\section{Dynamic Transport, Joint Contracts, and Priced Implementation}
\label{sec:r11_proofs}
This section supplies the proofs and numerical specifications for the section entitled ``Dynamic Adjustment and the Price of Commitment'' in the main paper. Main-paper labels in this independently compiled supplement are identified by their proposition titles; the substantive definitions are repeated here to avoid dependence on cross-document references.

\subsection{Proof of the Reachable-State Transport Proposition}
Let $V^\uparrow_n$ be the value with nonnegative deliberate drift, including the same elective stop action as the unrestricted problem. For a negative action $a$, let $a^0$ be its feasible zero-drift counterpart. Given enclosures $\ell\leq V^\uparrow_{n+1}\leq h$, write $V^\uparrow_{n+1}=v+z$, where $v=(h+\ell)/2$ and $|z|\leq e=(h-\ell)/2$. Entrywise positivity gives
\[
 (K^a-K^{a^0})V^\uparrow_{n+1}
 \leq (K^a-K^{a^0})v+|K^a-K^{a^0}|e.
\]
For each endpoint regime, maximize the reward difference over the two benefit endpoints. Taking the convex combination of these endpoint upper bounds bounds the difference under every mixture law. Maximizing that affine upper expression over the two law endpoints is valid. This is slightly more conservative than evaluating the absolute value of the mixed signed kernel itself, since
$|(1-\lambda)D_0+\lambda D_1|\leq(1-\lambda)|D_0|+\lambda|D_1|$.
The charged evaluation allowance then yields $Q^a(V^\uparrow_{n+1})<Q^{a^0}(V^\uparrow_{n+1})$.

At terminal time the unrestricted and restricted values agree. Suppose they agree at date $n+1$ on the common reachable set. Closure of the set under all feasible live transitions permits replacement of unrestricted continuation by restricted continuation in every date-$n$ backup. Every negative-drift backup is strictly dominated by a feasible zero-drift backup. All remaining operating and stop actions are available in the upward problem. Hence the two maxima agree. Backward induction proves equality on every reachable set, including either initial financial class: the counterpart has the same risky share, so it respects the class restriction. Boundary nodes have the common discharge value and require no strict operating-action inequality. The argument uses all feasible successors, not only successors of the ultimately chosen policy.

For completeness, the continuation enclosures used in the computation follow from a separate backward comparison. Start with the common terminal array. At each later date, evaluate the lower continuation at the lower benefit and higher fee and minimize over the law endpoints before maximizing over upward-feasible actions. This is a lower bound, because it permits an adverse law choice for each action. Evaluate the upper continuation at the higher benefit and lower fee and maximize over both law endpoints and actions. This is an upper bound. Add outward error padding after each stage. Neither recursion assumes that the optimal policy is constant on the box.

At the first date, the restricted objective is piecewise affine in the risky share for each fixed consumption--drift pair. The difference between the negative-drift action and its counterpart is piecewise affine on the union of their breakpoint sets. Its maximum over every admissible permission interval is therefore covered by the union of these knots and the largest-permission endpoints. A separately proposed action needs its own counterpart or direct domination proof. The deposited frozen proposals all have nonnegative drift on the reachable live states; this statement is checked rather than assumed.

\subsection{Continuation Incentives and the Dynamic Class Differential}
In a branch on which the exit fraction $\alpha$ is independent of deliberate preference drift, the arrival preference changes by $\alpha h\theta$, while consumption and the financial transition are fixed. The fundamental theorem of calculus for absolutely continuous, piecewise differentiable continuation gives the transport integral in the main paper. Integrating the quadratic flow cost gives $k\theta^2\E[(1-e^{-\rho\alpha h})/\rho]/2$. Immediate preference discharge requires a different branch calculation and is not covered by this smooth-transport interpretation; the signed-row test still applies.

An exact decomposition clarifies the dynamic comparison. Suppose the selected first financial controls $(c_\sigma,\pi_\sigma)$ agree between the adjusted and unadjusted class optima, as checked at the reported central contracts. Write $a_\sigma=(c_\sigma,\theta_\sigma,\pi_\sigma)$ and $a^0_\sigma=(c_\sigma,0,\pi_\sigma)$. Then
\begin{align}
 W^{\mathrm{adj}}_\sigma-W^0_\sigma
 ={}&Q^{a_\sigma}(V^{\mathrm{adj}}_1)
      -Q^{a^0_\sigma}(V^{\mathrm{adj}}_1)
      +K^{a^0_\sigma}(V^{\mathrm{adj}}_1-V^0_1).
 \label{eq:r11_s_decomp}
\end{align}
The last term measures exposure to future adjustment opportunities under the zero-drift initial action. Subtracting the two classes yields the reported relative local and future terms. If the financial controls differ, the additional term
$Q^{a^0_\sigma}(V^0_1)-W^0_\sigma$ must be retained; it is not generally zero. The replication code fails rather than silently suppressing that term when its equality premise fails.

For the boundary-state diagnostic, continuation is further written as $V_1=G+(V_1-G)$. This separates the liquidation benchmark from the liability of remaining in operation. The code independently reconstructs noncost flow, adjustment cost, immediate discharge, live liquidation-benchmark exposure, and continuation excess. The lower-preference-boundary mass is obtained by evaluating the boundary indicator through the stipulated first-date settlement lottery. It is discounted incidence, not an undiscounted probability and not a continuous-time hitting probability. The large full-minus-upward value difference is a reoptimized all-date comparison; the one-step component is not substituted for that full difference.

The primitive two-stage result and its common-cardinal-tilt robustness remain separate theorems. No step of this proof substitutes their cost coefficients or adjustment directions for those of the eight-date economy.

\subsection{Proof of the Joint-Region Certificate}
Fix a deterministic feasible stopping policy for the dates after the first, and a feasible first-date action. Affine dependence of rewards and kernels on the repeatedly drawn regime gives a degree-at-most-eight Bernstein polynomial in $\lambda$ by backward conditioning on the number of regime-1 draws. Conditional on the policy, the coefficients are affine in $d$ and $F$. The first-date coefficient is obtained by applying its selected first-date kernel to the seven-period continuation coefficients. Early stopping shortens effective degree but does not alter the degree-elevated representation.

At every benefit--fee corner, compute endpoint optima for the largest permissions. The signed corrected chord, or the count-information recursion with the same information timing, bounds each class value over the law interval. Convexity of optimized value in $(d,F)$ then bounds an interior reward point by the multilinear interpolation of these corner upper values. Both methods allow the same elective stop action at the same nodes. The first-date upper backup optimizes over the breakpoint menu corresponding to the largest permissions; the later upper backups use exactly the original full action menu.

For lower values, retain feasible policies optimized at the smallest permissions. They remain feasible at every permission pair in the box. Fix one such policy across all reward corners before taking a regional bound. Let $L^{p}_{c,j}$ be its restricted Bernstein coefficient at reward corner $c$ and law coefficient $j$, and let $U^\sigma_{c,j}$ be an upper coefficient for a comparison class. Nonnegativity and unit sum of the tensor interpolation weights imply
\[
 J_p-W_\sigma\geq\min_{c,j}\{L^p_{c,j}-U^\sigma_{c,j}\}.
\]
Taking the maximum of this bound over feasible lower policies is valid. Reversing the roles supplies an upper difference bound. Importantly, taking a separate best lower policy at each coefficient or reward corner before the minimum is not justified and is not done. The policy identity is held fixed through the complete corner-and-coefficient minimization.

For the adjusted and zero-adjustment spreads, apply this construction within each regime. For active surrender, apply it to the positive zero-adjustment class with surrender and to that same class with all eight intervals compulsory. The first-date permission extremes are treated in the same conservative manner. The resulting bounds include $2\times10^{-7}$ for each difference. Their displayed decimal bounds in the main paper are rounded outward.

If a policy has zero discounted elective-surrender incidence, its operating behavior and payoff are feasible in the noncancellable economy. It cannot exceed the latter class supremum. A strictly positive lower bound on the value gain from permitting surrender therefore implies that any policy closer to the surrender value than that bound must have positive elective-surrender incidence. This conclusion remains valid for a class defined by a supremum over strictly positive shares. The closure at zero is used only in the upper search; every retained lower policy has a strictly positive initial share in class $+$.

The proof is a continuum argument, not an inference from the six interior validation contracts. Those direct reoptimizations test the implementation against the derived bounds. The transport certificate supplies the independent statewise directional statement. Neither a point replay nor a class-value bound alone would supply it.

\subsection{Proof of the Priced Implementation Proposition}
For regime $r$, let $V^r_n$ denote the noncancellable continuation and let $\mathcal R_n$ be the feasible-support tube. A fee $F$ implements equality of surrender and noncancellable values on all eligible reachable nodes precisely when
\[
 F\geq F_r^*(m):=\max_{m\leq n<8}\max_{x\in\mathcal R_n\text{ live}}[G(x)-V^r_n(x)]_+.
\]
The empty maximum at $m=8$ is zero. This is a statewise threshold, not a necessary threshold for matching only the initial value. The earlier proof by backward induction applies without change to the enlarged first-date menu because the reachable tube now contains its successors. Policy inclusion gives $V^{\mathrm{adj}}_n\geq V^0_n$ and hence $F_{\mathrm{adj}}^*(m)\leq F_0^*(m)$.

Conditional on implementing the selected noncancellable allocation, $W^r$ and delivered service $A_r$ do not depend on the chosen instrument. For a given term, nondecreasing capacity cost makes the smallest sufficient capacity optimal. Participation requires
$q\geq G(x_0)-W^r+C(E)$ and $q\geq0$; hence the minimum grant is the stated positive part. The principal pays that grant and direct term cost. Minimizing over the eight terms proves the procurement formula, and comparison with $bA_r$ proves the break-even condition. At the threshold, continued operation must be the specified best-response selection. Requiring strict stopping incentives instead gives the infimum as fees decrease to the threshold; the code separately checks strictly higher fees and reports them as implementation tests, not as the exact continuous optimum.

When $z=G(x_0)-W^r\geq0$, the grant increment from capacity is exactly $C(E)$. For arbitrary $z$, monotonicity and unit Lipschitz continuity of $[\cdot]_+$ give
$0\leq[z+C(E)]_+-[z]_+\leq C(E)$. Comparing this increment with saved direct term cost proves the general instrument comparison, including the nonbinding-participation case. With $D(m)=\kappa(m-1)/8$ and binding participation, each term's cost is affine in $\kappa$. The lower envelope selects the economically preferred institution; crossings of these affine functions, not the free-term omission, determine its changes.

The illustrative coefficient $0.02$ and the values of $\kappa$ are theoretical primitives, not estimated enforcement costs. The principal's fees are zero on the implemented path, so no hidden fee revenue finances the grant. For active contracts outside this implementation family, any different fee recipient or principal revenue assumption must be entered explicitly.

\subsection{Arithmetic, Targets, and Reproduction}
The source fixes 1,617 continuation states, the inherited 1,568 later operating actions, eight dates, and the stored first-date breakpoint rows. The first rows are nonnegative with mass no larger than the inherited discount bound. The stopped reward includes the same boundary payoff and the elective fee. A full array manifest records data types, shapes, and SHA-256 identities for the primitive reward, duration, effort, CSR, and first-date arrays generated by the pinned sources.

The gamma-bound calculation in the R9 arithmetic audit is rerun at the new interval width and fee bound. Its per-class bound is smaller than the charged $10^{-7}$. The first-date backup has at most sixteen pre-coalescence transition entries per row. A signed difference of two such rows has at most thirty-two; its absolute-value transport bound is charged a further $2\times10^{-7}$ per comparison. The interval recursion is padded by $10^{-7}$ at each date. These pads exceed the round-to-nearest operation bound for the stored rewards and bounded eight-step continuation values. Positive interpolation weights and maxima do not amplify a uniform one-sided value error beyond the stated row-mass recursion. The support test uses the nonzero entries of these stored nonnegative kernels.

There is an additional distinction at the first date. The certified operator is the piecewise-affine interpolation of its deposited knot rows for each finite consumption--drift pair. Direct evaluation of the original transition formula at selected caps is checked against this representation; agreement is an implementation test, not an enclosure of the real-arithmetic constructor. Rebuilding with another environment requires checking the recorded array identities or treating the rebuilt arrays as a separately identified target. No constructor or diffusion enclosure is inferred from last-bit agreement.

The supplied runner executes the nine directional replays, the alternative-state witness, the dynamic decomposition, all twelve joint fee--permission interventions, both regional upper methods, the full reachable-state transport test, all eight enforcement alternatives, selected-control reconstructions, and interior-region validations. It saves the actual coefficient arrays and lower policies, not only their final extrema. The common endpoint bill, construction, and validation are charged to each matched upper-method workload. Kernel payload is reported separately from process peak memory, which is not measured. Historical training and full historical timing sweeps are preserved but are not described as newly rerun.

The proof establishes the indicated finite-array economic result. Transferring initial value inequalities to another target requires the additional value-error account stated in the main paper. Transferring directional dominance requires a compatible statewise transport and reachability account. These are different mathematical obligations.
